\section*{Abstract}
Retrieval-augmented generation (RAG) is increasingly used to answer questions from official guidance, where an answer often holds only under a condition. We ask where, in a chunked RAG pipeline, such conditions are lost. On ConditionalQA, turned into open retrieval over all \nPages{} of its gov.uk pages, the benchmark's human annotation lets us measure chunking, retrieval and evidence-reducing operators directly, and generation with the official metric under one shared prompt; an analysis plan fixed before any evaluation run defines the tests. Chunking separates the conditions of \sepCond\% of conditional questions from their answers, and losing the chunks that hold only a condition costs two local generators about five points of F1 with conditions. Retrieval is not biased against conditions, and the distractors it retrieves dilute the context more than its misses do. A reranker, a contradiction filter and a ConflictRAG-style resolver delete condition chunks as readily as answer chunks, yet change end-to-end scores little, because generation is the dominant loss: with every condition in context, even a 550B-parameter model leaves two thirds of them unstated. On RAMDocs, where discarding is correct, the same operators fail to isolate misinformation. A trained qualifier-aware context completion does not help, and its scorer does not transfer. We release all code, splits and model outputs.

\par\vspace{5pt}\noindent
\textbf{\fontsize{12}{12}\selectfont Keywords:} retrieval-augmented generation; conditional question answering; chunking; evidence selection; knowledge conflicts; evaluation
